\section{逆否、反证与分类讨论}
\label{sec:01-Mathematical-Language/03-Proof-Techniques/02}

\subsection{当你撞上一堵代数墙}

试试证明这个命题：``若整数 \(n^2\) 是偶数，则 \(n\) 是偶数。''

你自然地按直接证明的思路开局：假设 \(n^2\) 是偶数，即存在整数 \(k\) 使得
\[
n^2 = 2k.
\]
现在你要推出 \(n\) 是偶数，也就是存在某个 \(m\) 使得 \(n=2m\)。联系这两端的只能是 \(n=\sqrt{2k}\)。可是从平方根里扯出一个显式的整数因子 2？这条路走着走着就陷进去了——\(\sqrt{2k}\) 在整数范围内根本不闭合，你的代数工具箱在这里基本失灵。

不是你的技术不行，是你选错了入口。命题的逻辑形状决定了你该从哪扇门进去。有时候正面撞不开，绕到背面反而一推就开了。

\subsection{逆否命题：掉个头，同样真}

在逻辑上，\(P\to Q\) 和它的逆否命题 \(\neg Q\to\neg P\) 是完全等价的。两句话说的是一回事，只是语序颠倒了一下。这意味着：要证 \(P\to Q\)，去证 \(\neg Q\to\neg P\) 也一样。

回到那道题。原命题：\(n^2\) 偶 \(\to\) \(n\) 偶。逆否版本：\(n\) 不是偶数（即奇数）\(\to\) \(n^2\) 不是偶数（即奇数）。换句话说，就是证``奇数的平方还是奇数''。

这扇门就好进多了。设 \(n=2m+1\)，直接展开：
\[
n^2 = (2m+1)^2 = 4m^2 + 4m + 1 = 2(2m^2+2m) + 1.
\]
\(2m^2+2m\) 是整数，所以 \(n^2\) 确实是一个奇数。证完。

从直接路线的泥沼到逆否路线的坦途，分水岭只有一件事：不再试图从 \(n^2=2k\) 里往外扒拉东西，而是从 \(n\) 的表达式出发，自然地生成 \(n^2\) 的表达式。逆否法没有改变命题的真假，只是换了一个定义更容易上手的方向。别把它跟证逆命题 \(Q\to P\) 搞混了——那是完全不同的事情。

\subsection{反证法：假设结论是假的，然后让它活不下去}

再来一个更顽固的例子。证明：\(\sqrt{2}\) 不是有理数。

你无法直接证``不存在两个整数之比等于 \(\sqrt{2}\)''——存在命题可以举例子，不存在命题怎么举？你没办法枚举所有分数来逐个排除，那是无穷无尽的工作。直接路线再次被封死。

那就反过来。假设它\textbf{就是}有理数。那存在互素的整数 \(p,q\)（\(q\ne0\)），使得
\[
\sqrt{2} = \frac{p}{q}.
\]
这里``互素''——\(\gcd(p,q)=1\)——不是可有可无的细节。任何有理数都能约分到最简，我们只是明确地把这件事写下来。没有它，后面的矛盾就站不住脚。

两边平方：\(p^2 = 2q^2\)。所以 \(p^2\) 是偶数。用上一节的结论，\(p\) 本身也是偶数，设 \(p=2r\)。代回去：
\[
4r^2 = 2q^2 \quad\Longrightarrow\quad q^2 = 2r^2.
\]
那 \(q^2\) 也是偶数，\(q\) 也得是偶数。于是 \(p,q\) 都是偶数——但它们本该互素。两个数不可能同时都是偶数又互素。撞墙了。

这个矛盾的来源清晰而精确：从``\(\sqrt{2}\) 是有理数''这个假设出发，经过一系列无可挑剔的推导，撞上了``\(p,q\) 互素''这条你自己写下的前提。既然 \(\sqrt{2}\) 是有理数的假设会导致不可能的事，那这个假设本身就站不住脚。所以 \(\sqrt{2}\) 不是有理数。

反证法的力量来自一个明确的、无路可退的矛盾。你得能指着两行说：这两行不能同时成立。如果矛盾藏在模糊地带，那说明你的前提写得还不够干净。

\subsection{逆否法与反证法的分岔路口}

两者都是对付 \(P\to Q\) 的武器，但它们启动的前提不一样：

\begin{itemize}
\tightlist
\item
  逆否法：假设 \(\neg Q\)，目标是推出 \(\neg P\)。你手上只有``结论不成立''这一个假设。
\item
  反证法：假设 \(P\land\neg Q\)，目标是推出随便什么矛盾——可能是 \(R\land\neg R\)，可能是跟已知定理冲突，可能是跟定义冲撞。你手上的牌多了一张 \(P\)。
\end{itemize}

在经典逻辑里两者都成立，有时候写出来的证明还挺像。但有一条实用的分水岭：如果你能一路从 \(\neg Q\) 干净地推出 \(\neg P\)，逆否法更利落，不需要多余的假设；如果关键是 \(P\) 和 \(\neg Q\) 的相互作用制造了矛盾——就像 \(\sqrt{2}\) 的证明里，\(p^2=2q^2\) 这个等式来自前件（它是有理数），而互素的前提和后件的否定纠缠在一起——那反证法更自然。

\subsection{分类讨论：把一片混沌切成整齐的格子}

有些命题面对的论域太大了，一次处理不了。但如果你能把所有可能的情况切成几块，每一块内部定义变得简单、推导变得直白，那就分开来证。

比如证：对任意整数 \(n\)，\(n(n+1)\) 是偶数。

所有整数天然的就被分成偶数和奇数两类——既互斥又穷尽。那就按这个分类走：
\begin{itemize}
\tightlist
\item
  如果 \(n\) 是偶数，即 \(n=2k\)，那 \(n(n+1)=2k(2k+1)\)，成对地含有一个因子 2，偶数。
\item
  如果 \(n\) 是奇数，即 \(n=2k+1\)，那 \(n+1=2k+2=2(k+1)\)，同样含有一个因子 2，偶数。
\end{itemize}
两条路都通向``乘积是偶数''。所有整数都被覆盖了，没有漏网之鱼。

分类讨论的硬性要求是\textbf{穷尽}——漏掉一类，证明就有洞。互斥倒不是强制要求，重叠了顶多同一个对象被证了两遍，结论不会错。但好的分类通常还是尽量互斥，干净利落。

背后支撑按余数分类的是除法算法：每个整数除以正整数 \(m\)，商和余数唯一确定。所以余数 0, 1, \dots, \(m-1\) 恰好构成一个既不漏也不重的划分。按模 2 分对应奇偶性，按模 10 分对应末位数，按正负分对应符号相关的问题——分类的维度要贴着结论的需求来选。合适的分类会让每一块里的核心定义自动简化，剩下的推导往往就只是几步代数。

\subsection{不要把``没想到反例''当成证明}

反证法有一条严格的纪律，很多人第一次用的时候会犯：必须从否定假设逻辑地推出矛盾，而不是``看着就不像真的''。

``假设结论不成立，那显然不可能啊''——这不是反证，这是把结论换了个语气重说了一遍。``我们试了好多数都不行''——这也不是矛盾，这是数值实验的疲劳。真正的反证，得能指着两个具体陈述说：它们在逻辑上不能同时成立。比如``\(p,q\) 互素''和``\(p,q\) 都是偶数''。缺了这个精确的冲突点，你写出来的只是一个带感叹号的猜测。

换句话说，反证法不会因为你的信念强度而生效。它只看你能不能从否定假设出发，走向一个清晰、可验证、无法同时满足的逻辑死胡同。找不到那条死胡同，就别急着落笔写``矛盾''——先回头看看你的假设写全了没有。
